\documentclass[a4paper,11pt,fleqn]{article}%fleqn pour aligner les equations à gauche
\usepackage{graphicx}
\usepackage[utf8]{inputenc}  
\usepackage[T1]{fontenc}
\usepackage{lscape}

\usepackage{tgheros,textcomp}%Police Arial
\renewcommand{\familydefault}{\sfdefault}

%\setlength{\mathindent}{0pt}%indentation dans math = 0

\usepackage[top = 1cm, bottom = 2cm, left = 1cm, right = 1cm]{geometry}
\usepackage{titlesec}
%\setcounter{secnumdepth}{4}
\usepackage{xcolor}
\usepackage{amssymb} %double flèche de réaction
\usepackage{mathrsfs}%farraday
\usepackage{xtab}%tableau sur plusieurs pages
\usepackage{esvect}%grands vecteurs

\usepackage[version=4]{mhchem} %equations chimiques

\usepackage{array}
\usepackage{chemfig}
\usepackage{multirow}
%\usepackage{sectsty}

\usepackage{caption}
\usepackage{adjustbox}
\usepackage[labelformat=simple]{subfig}
\usepackage{float}


\usepackage{amsmath}
\usepackage{amssymb}
\usepackage[cal=boondoxo]{mathalfa} 
\usepackage{enumitem}
\usepackage{lastpage}
\usepackage{tcolorbox}
\usepackage{siunitx}
\usepackage[european, straightvoltages, RPvoltages]{circuitikz}
\usepackage{tikz}
\usetikzlibrary{shapes.geometric}
\usetikzlibrary{calc}
\usetikzlibrary{automata,positioning}
\usepackage{multicol}
\usepackage{eqnarray}
\usepackage{tabularx,colortbl}%colorer fond cellules

\sisetup{
	detect-all,
	output-decimal-marker={,},
	group-minimum-digits = 3,
	group-separator={~},
	number-unit-separator={~},
	inter-unit-product={~}
} %setup pour SIunitx

\usepackage{listings}%pour insérer le code python
\definecolor{codegreen}{rgb}{0,0.6,0}
\definecolor{codegray}{rgb}{0.5,0.5,0.5}
\definecolor{codepurple}{rgb}{0.58,0,0.82}
\definecolor{backcolour}{rgb}{0.95,0.95,0.92}


%\graphicspath{ {/home/rweye/Documents/Enseignement/2023/Cours/TSPE/Eval/Ressources/} }

\titleformat{\section}
{\titlerule
\vspace{.8ex}%
\normalfont \bfseries}
{\thesection :}{.5em}{\bfseries}

\renewcommand{\thesection}{Exercice \arabic{section}}
\titlespacing{\section}{0pt}{1em}{1em}
\renewcommand{\thesubsection}{\hspace{1cm}\alph{subsection}.}
\renewcommand\thesubfigure{\textbf{Document \arabic{subfigure} :}}

\title{\vspace{-1cm}\large 
	\begin{tabular}{|m{5cm}|>{\centering}m{8cm}|>{\centering\arraybackslash}m{4cm}|}
		\hline
		Nom : & \textbf{\textsc{DS TSPE N°1}} &  Note : \\
		Prénom : & & \\\hline
\end{tabular}}
\date{}

\newpagestyle{mystyle}{%
	\setfoot{ROSALIE}{\thepage\,$|$\,\pageref{LastPage}}{}
}
\renewpagestyle{plain}{%
	\setfoot{ROSALIE}{\thepage\,$|$\,\pageref{LastPage}}{}
}

\pagestyle{mystyle}

\newenvironment{itemize*}%
{\begin{itemize}[label=$-$]%
		\setlength{\itemsep}{0pt}%
		\setlength{\parskip}{0pt}}%
	{\end{itemize}}

\newenvironment{itemize**}%
{\begin{itemize}[label=,labelindent=0pt, leftmargin=*,topsep=0pt]%
		\setlength{\itemsep}{0pt}%
		\setlength{\parskip}{0pt}}%
	{\end{itemize}}

\renewcommand{\arraystretch}{1.5}


\begin{document}
\maketitle
\vspace{-2.5cm}

\section*{Exercice : Cuisine et acidité}
Le chou rouge est un légume riche en fibres et en vitamines, qui se consomme en salade ou cuit. Mais la cuisson du chou rouge peut réserver des surprises : chou rouge et eau de cuisson deviennent bleus. Pour rendre au chou sa couleur violette, on peut ajouter un filet de jus de citron ou du vinaigre.\\
Une autre modification de couleur peut surprendre le cuisinier : versée dans un évier contenant du détergent, l'eau de cuisson du chou devient verte.
	\subsection{Acides et bases dans la cuisine}
	\begin{enumerate}
		\item Préciser la propriété principale d'un indicateur coloré acide-base.
		\item Justifier le caractère basique du détergent.
		\item Indiquer le schéma de Lewis des deux espèces du couple acide base de l'acide éthanoïque \ce{CH3COOH}.
	\end{enumerate}
	
	\subsection{Dilution du vinaigre}
	Le vinaigre commercial à \SI{1.0}{\mole\per\liter} est trop concentré pour être étudié. On le dilue donc dix fois. On dispose de la verrerie suivante :
	\begin{itemize}
		\item éprouvettes (mL) : 5 ; 10 ;25 ; 50 ; 100
		\item pipettes jaugées (mL) : 1,0 ; 5,0 ; 10,0 ; 20,0
		\item fioles jaugées (mL) : 150,0 ; 200,0 ; 250,0 ; 500,0 
	\end{itemize}
	\begin{enumerate}[resume]
		\item Proposer un protocole de dilution à partir de cette liste la verrerie.
	\end{enumerate}
	
	\subsection{Réaction du vinaigre sur la soude}
	On ajoute un volume $V_{eau}$ = \SI{60}{mL} afin d'immerger les électrodes du pH-mètre à un volume $V_A$ = \SI{10,0}{mL} de la solution diluée de vinaigre.\\
	Tout en mesurant le pH, on ajoute progressivement une solution aqueuse d'hydroxyde de sodium (\ce{Na+(aq)} ; \ce{HO-(aq)}) de concentration en soluté apporté $c$ = \SI{1.0E-1}{\mole\per\liter}. On a alors réaction entre les ions hydroxyde et l'acide éthanoïque contenu dans le vinaigre.\\
	La courbe $pH = f(V)$ du document 2 est alors obtenue.
	\begin{enumerate}[resume]
		\item Ecrire  l'équation associée à la réaction.
		\item Pour $V$ = \SI{6.0}{mL}, déterminer l'avancement maximal $x_{max}$.
		\item Préciser si la réaction étudiée est totale pour un volume versé $V$.
		\item Déterminer théoriquement le volume équivalent à partir de l'énoncé.
		\item Le groupe d'élève a t il réalisé correctement le dosage ? Proposer des pistes d'amélioration pour leur méthode.
	\end{enumerate}
	
	\subsection{Titre massique du vinaigre}
	\begin{enumerate}[resume]
		\item  A partir des données, calculer le titre massique du vinaigre en degré (\unit{\gram\per 100\gram})
	\end{enumerate}

\noindent \colorbox{red!10}{\begin{minipage}{.97\textwidth}
\textbf{Données :}
\begin{itemize}
	%\item \textbf{Relation utile :} $\dfrac{([\ce{H3O+}]\cdot[\ce{HO-}])}{{c^0}^2} = 10^{-14}$
	\item \textbf{Masses molaires atomiques : } $M(C)$ = \SI{12.0}{\gram\per\mole}, $M(H)$ = \SI{1.0}{\gram\per\mole}, $M(O)$ = \SI{16.0}{\gram\per\mole}
	\item \textbf{Densité du vinaigre : } $d$ = \SI{1.01}{\kilo\gram\per\deci\meter\cubed}
\end{itemize}
\end{minipage}
}
\vspace{1em}

\noindent \colorbox{orange!20}{\begin{minipage}{.47\textwidth}
\textbf{Document 1 : Chou rouge}\\
La première utilisation d'un indicateur coloré pour les titrages acide-base a été faite en 1767 par W. Lewis. il employait un extrait de tournesol. On utilisait à l'époque des extraits de plantes qui changent de couleur avec l'acidité du milieu.\\
En voici quelques-uns parmi les plus connus et les meilleurs : l'artichaut, la betterave rouge et la chou rouge, de loin l'extrait le plus intéressant car sa couleur change nettement suivant la valeur du pH.\\
\begin{tabular}{|>{\centering}m{1.5cm}|>{\centering}m{1cm}|>{\centering}m{1cm}|>{\centering}m{1cm}|>{\centering}m{1cm}|>{\centering\arraybackslash}m{1cm}|} \hline
	\textbf{pH} & 0-3 & 4-6 & 7-8 & 9-12 & 13-14\\ \hline
	\textbf{Couleur} & rouge & violet & bleu & vert & jaune \\ \hline
\end{tabular}
\end{minipage}
}
\colorbox{green!20}{\begin{minipage}{.47\textwidth}
\textbf{Document 2 : Suivi pH-métrique d'un groupe d'élèves}
\begin{center}
	\includegraphics[scale=.5]{Ressources/TSPE_DS1c_Fig1.png}
\end{center}
\end{minipage}
}
	
\end{document}